% Figure: the case-study parabolic suspension cable. A cable hangs between
% two level towers a span L apart, carrying a uniform deck load w per unit
% horizontal length, sagging by h at midspan. The horizontal tension H is
% constant along the cable; the cable tension is largest at the towers,
% where the cable is steepest. The free body of the half-cable shows H at
% midspan, the vertical reaction wL/2 at the tower, and the tower tension T.
\begin{figure}[htbp]
\centering
\begin{tikzpicture}[>=latex, x=0.85cm, y=0.85cm]
  % towers at x=0 and x=12, deck datum at y=0
  % parabola: y = -h*(1 - (1 - 2x/L)^2) shifted; draw sag below the chord
  % use L=12 units, sag h=2.4 units (drawn). cable low point at midspan.
  \def\Lspan{12}
  \def\sag{2.4}
  % tower tops
  \coordinate (TL) at (0,0);
  \coordinate (TR) at (\Lspan,0);
  % parabolic cable, sampled (y dips by sag at midspan)
  \draw[engblue, very thick]
    plot[domain=0:12, samples=49]
      (\x, {-\sag*(1 - (1 - 2*\x/12)*(1 - 2*\x/12))});
  % towers
  \draw[engblack, line width=1.3pt] (0,0) -- (0,-3.2);
  \draw[engblack, line width=1.3pt] (\Lspan,0) -- (\Lspan,-3.2);
  % deck (the suspended load datum), drawn as a light line under the cable
  \draw[enggray, thick] (0,-3.2) -- (\Lspan,-3.2);
  % hangers (vertical ties from cable to deck), a few representative
  \foreach \x in {1.5,3,4.5,6,7.5,9,10.5} {
    \draw[enggray, thin]
      (\x, {-\sag*(1 - (1 - 2*\x/12)*(1 - 2*\x/12))}) -- (\x,-3.2);
  }
  % distributed deck load arrows
  \foreach \x in {1,2,...,11} {
    \draw[->, engred, thin] (\x,-3.95) -- (\x,-3.25);
  }
  \node[engred, font=\scriptsize, anchor=north] at (6,-4.0) {$w$ per unit horizontal length};
  % low point at midspan
  \coordinate (LP) at (6,-\sag);
  \fill[engblack] (LP) circle (1.6pt);
  % horizontal tension H at the low point (acts along the cable, horizontal)
  \draw[->, very thick, englime!70!black] (LP) -- ++(-1.6,0);
  \node[englime!60!black, font=\scriptsize, anchor=north] at ($(LP)+(-1.0,-0.05)$) {$H$};
  % tower tension T tangent at the left support
  \draw[->, very thick, engamber] (TL) -- ++(-1.2,0.95);
  \node[engamber, font=\scriptsize, anchor=south east] at ($(TL)+(-0.7,0.5)$) {$T_{\max}$};
  % support reaction (vertical) at left tower
  \draw[->, thick, englime!70!black] (TL) -- ++(0,1.1);
  \node[font=\scriptsize, anchor=west] at ($(TL)+(0.05,0.7)$) {$V=\tfrac{wL}{2}$};
  % joints
  \fill[engblack] (TL) circle (1.6pt);
  \fill[engblack] (TR) circle (1.6pt);
  % sag dimension
  \draw[<->, enggray, thin] (12.6,0) -- (12.6,-\sag);
  \node[enggray, font=\scriptsize, anchor=west] at (12.6,{-\sag/2}) {sag $h$};
  % span dimension
  \draw[<->, enggray, thin] (0,1.6) -- (\Lspan,1.6);
  \node[enggray, font=\scriptsize, anchor=south] at (6,1.6) {span $L$};
\end{tikzpicture}
\caption[Parabolic suspension cable under deck load]{The parabolic
suspension cable of the case study. A cable hung between two level towers
a span $L$ apart carries a deck load $w$ per unit horizontal length through
closely spaced hangers, so the load on the cable is uniform in the
horizontal coordinate and the cable takes a parabola rather than the
catenary of a cable under its own weight. The horizontal component $H$ of
the cable tension is the same at every point, because no horizontal load
acts between the towers; the cable tension is therefore smallest at the
midspan low point, where it equals $H$, and largest at the towers, where
the added vertical component $wL/2$ makes the tension $T_{\max}=\sqrt{H^2 +
(wL/2)^2}$. Each tower carries a vertical reaction $wL/2$ and must also
resist the horizontal pull $H$, which a back-stay or an anchored foundation
supplies.}
\label{fig:vol03ch02:parabolic-cable}
\end{figure}
